\addtocontents{toc}{\protect\setcounter{tocdepth}{0}}
% \Needspace{4\baselineskip}
\subsection*{Literatūros apžvalgos apie giluminio neuroninio tinklo metodą akustinių anomalijų įvykiams aptikti}\label{chap:onelt}
% \addtocontents{toc}{\protect\setcounter{tocdepth}{2}}
Akustinių metodų naudojimas aplinkos stebėjimo ar pažeidimų nustatymo kontekste turi reikšmingą istorinį precedentą~\cite{Simonovic2021, Vicuna2017}. Paprastai neįprastos sąlygos stebimose įrangose gali būti nustatytos pagal pagamintų akustinių signalų savybių kitimą, pavyzdžiui, dažnio ir amplitudės pokyčius.~\cite{Holford2017, Cooper2020}. Akustinio metodo pranašumas, palyginti su kitais metodais, yra tai, kad \dots

\dots
